% Generated draft: requires mathematical and source review.
\documentclass[11pt]{article}
\usepackage[margin=25mm]{geometry}
\usepackage{fontspec}
\setmainfont{Latin Modern Roman}
\usepackage{amsmath,amssymb,mathtools,unicode-math}
\setmathfont{Latin Modern Math}
\usepackage{enumitem,needspace,xurl,hyperref}
\setlength{\parindent}{0pt}
\setlength{\parskip}{5pt}
\setlength{\emergencystretch}{3em}
\newcommand{\N}{\mathbb{N}}
\newcommand{\Z}{\mathbb{Z}}
\newcommand{\Q}{\mathbb{Q}}
\newcommand{\R}{\mathbb{R}}
\newcommand{\C}{\mathbb{C}}
\newcommand{\F}{\mathbb{F}}
\newcommand{\A}{\mathbb{A}}
\newcommand{\PP}{\mathbb{P}}
\newcommand{\E}{\mathbb{E}}
\newcommand{\eps}{\varepsilon}
\newcommand{\dd}{\,\mathrm{d}}
\begin{document}
\section*{UnsolvedMath ID: 800354}
\textbf{Source problem code:} catalogue entry
\subsection{Definitions and mathematical statement}
For each integer \(n\geq 1\), let \(R(n)\) denote the number described in the supplied source as the number of ``royal paths in a lattice of order \(n\),'' with the intended definition referenced to OEIS entry A006319. Retain \(T_{112}(a,b)\) as the notation used in the source on trimer covers in the triangular grid.

The source asks whether
\[
T_{112}(n+2,2n)=R(n)\qquad\text{for every integer }n\geq 1.
\]
The supplied record reports that this equality has been verified for \(1\leq n\leq 15\); that verification is a source claim and has not been independently checked here.

The supplied record does not define \(T_{112}(a,b)\), its underlying region or admissible covers, or the meaning of its subscript. It also does not define the royal paths, their lattice, permitted steps, endpoints, or boundary restrictions. Consequently, the displayed equality preserves the source question but is not a self-contained mathematical formulation until the intended definitions are supplied from the cited references. No choices for those missing definitions are imposed here.
\subsection{Short English statement}
The question attributed to Propp asks whether \(T_{112}(n+2,2n)\) counts the royal paths referenced by OEIS A006319 for every integer \(n\geq 1\). The record reports agreement for \(1\leq n\leq 15\) and labels the problem open. The definitions of both counting objects require clarification from the cited sources.
\subsection{Sources}
\begin{itemize}
\item Local UnsolvedMath registry: supplied complete record for permanent ID \(800354\), exact source code \texttt{catalogue entry}, titled ``Trimer Covers in the Triangular Grid --- Problem 13.'' Attribution: Propp, proposed year 2022. The registry labels the problem open; this supplied status is not an independently verified literature assessment.
\item Propp, \emph{Trimer Covers in the Triangular Grid: Twenty Mostly Open Problems} (2022), source item ``Problem 13 (\#6).'' Supplied source URL: \url{https://arxiv.org/abs/2206.06472}. The record reports access on 2026-10-16 and extraction by \texttt{arxiv\_latex}; these are supplied provenance details, not a claim of access or verification here.
\item OEIS entry A006319, referenced in the supplied statement for ``royal paths in a lattice of order \(n\)'': \url{https://oeis.org/A006319}. The entry's contents were not included in the supplied record.
\end{itemize}
\end{document}
